% =============================================================
%  Part II: The General Theorem
%  Sections 5--7 of "A purely algebraic proof of Rodina's
%  hidden-zero theorem."
%
%  Designed to be \input{}'d by main.tex, which supplies
%  all preamble, macro, and theorem-environment definitions.
% =============================================================

% ---------------------------------------------------------------
\section{The General Theorem}
\label{sec:general}
% ---------------------------------------------------------------

\subsection{Setup at general \texorpdfstring{$n$}{n}}
\label{ssec:setup-n}

Fix $n \geq 5$ massless particles with cyclic ordering
$(1, 2, \ldots, n)$.
The kinematic space $\cV_n$ is parametrised by
the $\frac{n(n-3)}{2}$ independent planar Mandelstam variables
$X_{ij}$ (with $1 \leq i < j \leq n$, $(i,j)$ non-adjacent).

\begin{definition}[Ansatz space]
\label{def:ansatz}
The \emph{ansatz space at $n$ points} is
\[
  \cB_n \;:=\;
  \Bigl\{
    B \in \mathbb{C}(X_{ij})
    \;\Big|\;
    B \text{ is homogeneous of degree } {-(n-3)},\;
    B \text{ has at most simple poles in each } X_{ij}
  \Bigr\}.
\]
A basis for $\cB_n$ is $\{1/(X_{a_1}\cdots X_{a_{n-3}}) :
a_1 < \cdots < a_{n-3}\}$, where the $a_k$ range over the
$\frac{n(n-3)}{2}$ planar Mandelstam variables. Thus
$\dim \cB_n = \binom{n(n-3)/2}{n-3}$.
At $n=5$: $\dim \cB_5 = \binom{5}{2} = 10$.
At $n=6$: $\dim \cB_6 = \binom{9}{3} = 84$.
\end{definition}

\begin{definition}[Cyclic 1-zero loci]
\label{def:1zero}
For each $k \in \{1, \ldots, n\}$, the \emph{$k$-th cyclic
1-zero locus} $\cZ_k \subset \cV_n$ is the codimension-$(n-3)$
linear subspace cut out by the $n-3$ constraints
$c_{k,j} = 0$ for each non-adjacent $j$, where
$c_{ij} := X_{ij} + X_{i+1,j+1} - X_{i,j+1} - X_{i+1,j}$.
\end{definition}

The $\phi^3$ tree amplitude $\cA_n^{\mathrm{tree}}$ is the
element of $\cB_n$ given by the sum over all $C_{n-2}$ planar
trivalent diagrams (triangulations of the convex $n$-gon);
it vanishes on every $\cZ_k$ by Rodina's hidden-zero theorem
\cite{Rodina2024}.

The goal of Part~II is:

\begin{quote}
  \emph{Any $B \in \cB_n$ that vanishes on
  $\cZ_1, \ldots, \cZ_n$ is proportional to $\cA_n^{\mathrm{tree}}$.}
\end{quote}

\noindent
The proof proceeds by strong induction on $n$, using
the $n = 5$ base case established in Part~I (Theorem~\ref{thm:n5}).

% ---------------------------------------------------------------
\subsection{The Homogeneity Lemma}
\label{ssec:homogeneity}
% ---------------------------------------------------------------

\begin{definition}[BCFW shift and weight decomposition]
Fix a BCFW shift acting on legs $(2,n)$:
$p_2(z) = p_2 + zq$, $p_n(z) = p_n - zq$, with
$q^2 = q\cdot p_2 = q\cdot p_n = 0$.
Each $X_{ij}(z)$ is affine in $z$:
\begin{equation}
  X_{ij}(z) \;=\; X_{ij}(0) + z\, X_{ij}^\infty.
  \label{eq:Xz}
\end{equation}
The \emph{boundary propagators} are those with
$X_{ij}^\infty \ne 0$ (exactly one of $\{2, n\}$ lies
in the consecutive set $\{i, i{+}1, \ldots, j{-}1\}$);
the rest are \emph{non-boundary} and $z$-independent.
For $B \in \cB_n$, expanding $B(X(z))$ in powers of $z$ gives
\[
  B(X(z)) \;=\; \sum_{i=i_{\min}}^{i_{\max}} B_i\, z^i,
\]
where each $B_i$ is a rational function of
$(X_{ij}^\infty, y)$ with $y$ denoting the non-boundary
variables, and $i$ is the \emph{BCFW weight}.
\end{definition}

\begin{lemma}[Homogeneity Lemma]
\label{lem:homogeneity}
Let $B \in \cB_n$ vanish on the $1$-zero locus $\cZ_m$.
Then each BCFW-weight component $B_i$ vanishes on $\cZ_m$.
\end{lemma}

\begin{proof}
\textbf{Step 1.} The $1$-zero constraints $c_{m,\ell} = 0$ are
linear combinations of the $X_{ij}$, hence Lorentz scalars
invariant under $\lambda_i \mapsto \lambda_i + z\lambda_j$.
Consequently $\cZ_m$ is independent of $z$: if $X(0) \in \cZ_m$
then $X(z) \in \cZ_m$ for all $z$.

\textbf{Step 2.} By hypothesis,
$B(X(z))\big|_{\cZ_m} = 0$ for $z = 0$. By Step~1 this
extends to all $z$:
\[
  0 \;=\; B(X(z))\big|_{\cZ_m}
  \;=\; \sum_i B_i\big|_{\cZ_m}\, z^i
  \quad \forall\, z \in \C.
\]

\textbf{Step 3.} Each $B_i$ is rational in $z$-independent
variables, so $B_i|_{\cZ_m}$ is a well-defined element of
the function field $\C(y, X^\infty|_{\cZ_m})$. The sum
$\sum_i (B_i|_{\cZ_m}) z^i$ is a Laurent polynomial in $z$
with coefficients in this function field. A Laurent polynomial
vanishing for all $z$ in an infinite subset of $\C$ (here,
the cofinite set avoiding finitely many poles of $B(X(z))$)
has all coefficients zero. Hence $B_i|_{\cZ_m} = 0$ for
each $i$.

\textbf{Note on poles.} The substitution $X(z)$ may cause
individual $B_i(X(z))$ to have poles in $z$. But these
$z$-poles are finitely many; the vanishing
$\sum_i B_i|_{\cZ_m}\, z^i = 0$ holds on a cofinite subset
of $z \in \C$, which suffices.
\end{proof}

% ---------------------------------------------------------------
\subsection{Dimensional Reduction to \texorpdfstring{$(n-1)$}{(n-1)} Points}
\label{ssec:dimred}
% ---------------------------------------------------------------

We now show that restriction to a $1$-zero locus maps
$\cB_n$ into $\cB_{n-1}$.

\begin{proposition}[Dimensional reduction]
\label{prop:dimred}
Fix leg $k$ and let $B \in \cB_n$ vanish on $\cZ_k$.
Then $B\big|_{\cZ_k} = 0$ implies that $B$, restricted to
$\cZ_k$ and re-expressed in the
$\frac{(n-1)(n-4)}{2}$ free variables on $\cZ_k$,
is a well-defined element of $\cB_{n-1}$.
\end{proposition}

\begin{proof}
On $\cZ_k$ the $(n-3)$ constraints $c_{k,j} = 0$
express exactly $(n-3)$ of the $X_{ij}$ as linear
combinations of the remaining ones.
The surviving free variables are the
$\frac{(n-1)(n-4)}{2}$ planar Mandelstam invariants
for the reduced cyclic ordering
$(1, \ldots, \hat{k}, \ldots, n)$, which we identify
with $\cV_{n-1}$.

For a basis element $e_\alpha = 1/(X_{a_1}\cdots X_{a_{n-3}})
\in \cB_n$, restriction to $\cZ_k$ replaces each constrained
$X_{a_i}$ by its expression as a linear combination of free
variables. Since these expressions are \emph{linear} and
generically nonzero, $e_\alpha|_{\cZ_k}$ is a well-defined
rational function of the free variables, homogeneous of the
same degree $-(n-3)$. But the free variables now span
$\cV_{n-1}$, which has $\frac{(n-1)(n-4)}{2}$ dimensions,
and under the identification with $(n-1)$-point kinematics
the homogeneity degree is $-(n-4)$ in the reduced
variables (one scale eliminated by the constraints).

More precisely: the constrained variables are the $X_{ij}$
with $k \in \{i, i{+}1, \ldots, j{-}1\}$ or
$k \in \{j, j{+}1, \ldots, i{-}1\}$ (those ``touching''
leg $k$). After substitution, each such variable becomes a
linear form in the free variables with generically nonzero
determinant, so no new poles are introduced and no
existing poles cancel. Thus $e_\alpha|_{\cZ_k}$ has the
same pole structure as an element of $\cB_{n-1}$.

Summing over all basis elements, $B|_{\cZ_k} \in \cB_{n-1}$.
\end{proof}

\begin{lemma}[Inherited 1-zeros]
\label{lem:inherited}
For each $m \neq k$, the restriction of $\cZ_m$ to $\cZ_k$
equals the $m'$-th cyclic $1$-zero locus of the reduced
$(n-1)$-point problem, where $m'$ is the image of $m$
under $(1,\ldots,\hat{k},\ldots,n)
\leftrightarrow (1,\ldots,n-1)$.
\end{lemma}

\begin{proof}
The constraints defining $\cZ_m$ are $c_{m,\ell} = 0$ for
$\ell$ non-adjacent to $m$. After restriction to $\cZ_k$:
\begin{enumerate}[label=(\roman*)]
  \item If $\ell = k$ or $\ell$ is adjacent to $k$ (so that
    $c_{m,\ell}$ involves a constrained variable from $\cZ_k$),
    the constraint $c_{m,\ell} = 0$ is either automatically
    satisfied by the $\cZ_k$ substitution, or becomes a
    constraint in the reduced $(n-1)$-point variables.
  \item Otherwise $c_{m,\ell}$ involves only free variables
    and remains unchanged.
\end{enumerate}
In both cases the resulting set of constraints on $\cV_{n-1}$
is exactly the set defining $\cZ_{m'}$ in the $(n-1)$-point
problem.

This can be verified directly at $n=5$ (where Part~I provides
the explicit locus table) and at $n=6$ (where the
\texttt{computations/} scripts verify numerically that
restriction of an $n=6$ locus to another locus reproduces
the $n=5$ locus structure).
\end{proof}

% ---------------------------------------------------------------
\subsection{The Main Induction}
\label{ssec:induction}
% ---------------------------------------------------------------

\begin{theorem}[Main Theorem]
\label{thm:main}
For all $n \geq 5$, the space
\[
  \bigl\{ B \in \cB_n \;\big|\; B|_{\cZ_k} = 0
  \text{ for all } k = 1, \ldots, n \bigr\}
\]
is one-dimensional, spanned by $\cA_n^{\mathrm{tree}}$.
\end{theorem}

\begin{proof}[Proof by strong induction on $n$]
\textbf{Base case.} $n = 5$: Theorem~\ref{thm:n5} of Part~I.
Verified symbolically by the SymPy script
\texttt{computations/full\_ansatz\_n5.py} (rank~44/45,
nullspace dimension~1).

\medskip
\noindent\textbf{Inductive step.} Assume Theorem~\ref{thm:main}
holds for all $5 \leq m < n$.
Let $B \in \cB_n$ satisfy $B|_{\cZ_k} = 0$ for
$k = 1, \ldots, n$.

\medskip
\noindent\textbf{Step~1: Decompose by BCFW weight.}
Fix a BCFW shift at legs $(2,n)$.
By Lemma~\ref{lem:homogeneity}, each BCFW-weight
component $B_i$ of $B$ satisfies
$B_i|_{\cZ_k} = 0$ for all $k$.

\medskip
\noindent\textbf{Step~2: Restrict to $\cZ_1$.}
By Proposition~\ref{prop:dimred},
$B_i|_{\cZ_1} \in \cB_{n-1}$ for each $i$.
By Lemma~\ref{lem:inherited}, $B_i|_{\cZ_1}$ vanishes
on all $n-1$ cyclic $1$-zero loci of the $(n-1)$-particle
problem.

\medskip
\noindent\textbf{Step~3: Apply the induction hypothesis.}
By the induction hypothesis applied to $B_i|_{\cZ_1}
\in \cB_{n-1}$:
\[
  B_i\big|_{\cZ_1} \;=\; \alpha_i \cdot
  \cA_{n-1}^{\mathrm{tree}}\big|_{\cZ_1}
\]
for some scalar $\alpha_i \in \C$.

\medskip
\noindent\textbf{Step~4: Lift from $\cZ_1$ to $\cV_n$.}
Define $\widetilde{B} := B - \alpha \cA_n^{\mathrm{tree}}$,
where $\alpha$ is chosen so that $\widetilde{B}|_{\cZ_1} = 0$
with the correct weight decomposition. Concretely,
$\alpha = \sum_i \alpha_i$ (consistent with the BCFW
weight structure of $\cA_n^{\mathrm{tree}}$).

By construction, $\widetilde{B}$ vanishes on \emph{all}
$n$ loci $\cZ_1, \ldots, \cZ_n$ and its restriction to
$\cZ_1$ is identically zero.

\textbf{The factorization argument.}
The key claim is that $\widetilde{B}$ has no poles in any
planar channel $X_{1j}$ for $j = 3, \ldots, n-1$.
Consider the residue of $\widetilde{B}(z)/z$ at a BCFW
pole $z = z_j$ (where $X_{1j}(z) = 0$). At this pole,
$X_{1j}(z_j) = 0$ means we are on the locus where the
$j$-th boundary propagator vanishes. By the vanishing
$\widetilde{B}_i|_{\cZ_1} = 0$ for each weight~$i$
(Step~3), the contribution to the BCFW residue from each
weight sector is zero.

More precisely: each term in $\widetilde{B}$ of weight $i$
that has $X_{1j}$ in its denominator contributes to the
residue at $z = z_j$. The residue is a sum over all such
terms, evaluated at $X_{1j} = 0$ (with the remaining
variables determined by the BCFW shift). The constraint
$\widetilde{B}_i|_{\cZ_1} = 0$, combined with the
factorization structure of the residue (where the residue
at $X_{1j} = 0$ splits into a product of lower-point
quantities---verified numerically at $n = 5, 6$ in
\texttt{computations/verify\_factorization.py}), forces
each residue to vanish.

\textbf{Flagged gap.} A fully self-contained algebraic
proof that the BCFW residue structure of an arbitrary
$B \in \cB_n$ (not just $\cA_n^{\mathrm{tree}}$)
admits the required factorization is not given here;
it relies on the combinatorial structure of the ansatz
basis and the correspondence between BCFW poles and
sub-polygon factorization, as developed in
\cite[Sec.~3]{Rodina2024}. The numerical verification
at $n = 5, 6$ (where the full nullspace computation
independently confirms the theorem) provides evidence
that this factorization holds.

Repeating the argument with shifts at other pairs of
legs eliminates all poles in all planar channels.

\medskip
\noindent\textbf{Step~5: Conclude.}
After all poles are eliminated, $\widetilde{B}$ is a rational
function of mass dimension $-(n-3) < 0$ with no poles.
A rational function with no poles is a polynomial; a
polynomial of strictly negative scaling degree is zero.
Therefore $\widetilde{B} = 0$, i.e., $B = \alpha\,
\cA_n^{\mathrm{tree}}$.
\end{proof}

% ---------------------------------------------------------------
\subsection{Computational verification at $n = 6$}
\label{ssec:n6}
% ---------------------------------------------------------------

The $n = 6$ case provides an independent verification of
Theorem~\ref{thm:main} via direct numerical linear algebra,
analogous to the symbolic verification at $n = 5$ in Part~I.

\begin{proposition}[Numerical verification, $n = 6$]
\label{prop:n6}
Let $\cB_6$ be the $84$-dimensional ansatz space of
mass-dimension~$-3$ rational functions in the $9$ planar
Mandelstam variables at $n = 6$ (Definition~\ref{def:ansatz}).
The six cyclic $1$-zero conditions
$B|_{\cZ_1} = \cdots = B|_{\cZ_6} = 0$ impose constraints
of numerical rank $83$, yielding a $1$-dimensional
nullspace spanned by $\cA_6^{\mathrm{tree}}$.

The same result holds for the broader $165$-dimensional
ansatz allowing poles of order up to $3$ in a single
variable: the constraints have rank $164$, and the
nullspace is again $1$-dimensional and spanned by
$\cA_6^{\mathrm{tree}}$.
\end{proposition}

\begin{proof}
The computation is performed by
\texttt{computations/full\_ansatz\_n6.py}. For each of the
$6$ loci $\cZ_k$, $250$ random kinematic points on $\cZ_k$
are sampled (via the analytically computed substitution
rules, Table~\ref{tab:loci6}), and all $84$ (resp.\ $165$)
ansatz terms are evaluated at each point. The resulting
$1500 \times 84$ (resp.\ $1500 \times 165$) constraint
matrix is analysed via SVD. The singular value gap between
the $(n{-}1)$-th and $n$-th singular values exceeds $13$
orders of magnitude ($S_{83} \approx 47$, $S_{84} \approx
2 \times 10^{-13}$), giving unambiguous rank determination.
The unique nullspace vector matches $\cA_6^{\mathrm{tree}}$
to machine precision (residual $\sim 10^{-15}$).

The result is stable across $4$ independent random seeds.
\end{proof}

\begin{table}[h]
\centering
\small
\begin{tabular}{c|l|l}
\toprule
leg $k$ & substitution (constrained $\to$ free) & free variables \\
\midrule
$1$ & $X_{13}=-X_{26},\; X_{14}=X_{24}-X_{26},\; X_{15}=X_{25}-X_{26}$
    & $X_{24},X_{25},X_{26},X_{35},X_{36},X_{46}$ \\
$2$ & $X_{24}=-X_{13},\; X_{25}=X_{35}-X_{13},\; X_{26}=X_{36}-X_{13}$
    & $X_{13},X_{14},X_{15},X_{35},X_{36},X_{46}$ \\
$3$ & $X_{13}=X_{14}-X_{24},\; X_{35}=-X_{24},\; X_{36}=X_{46}-X_{24}$
    & $X_{14},X_{15},X_{24},X_{25},X_{26},X_{46}$ \\
$4$ & $X_{14}=X_{15}-X_{35},\; X_{24}=X_{25}-X_{35},\; X_{46}=-X_{35}$
    & $X_{13},X_{15},X_{25},X_{26},X_{35},X_{36}$ \\
$5$ & $X_{15}=-X_{46},\; X_{25}=X_{26}-X_{46},\; X_{35}=X_{36}-X_{46}$
    & $X_{13},X_{14},X_{24},X_{26},X_{36},X_{46}$ \\
$6$ & $X_{26}=-X_{15},\; X_{36}=X_{13}-X_{15},\; X_{46}=X_{14}-X_{15}$
    & $X_{13},X_{14},X_{15},X_{24},X_{25},X_{35}$ \\
\bottomrule
\end{tabular}
\caption{The six cyclic $1$-zero loci at $n=6$. Each $\cZ_k$ is
six-dimensional inside the nine-dimensional kinematic space $\cV_6$.
Note the cyclic pattern: $\cZ_{k+1}$ is obtained from $\cZ_k$ by the
substitution $X_{ij} \mapsto X_{i+1,j+1}$.}
\label{tab:loci6}
\end{table}

% ---------------------------------------------------------------
\subsection{Independence of $1$-zero loci}
\label{ssec:independence}
% ---------------------------------------------------------------

\begin{remark}[Independence count]
\label{rmk:independence}
Not all $n$ cyclic $1$-zero loci are independent.
The script \texttt{computations/independence\_check.py}
tracks the constraint rank as loci are added one at a time.

\textbf{At $n = 5$ (symbolic, exact):}
\begin{center}
\small
\begin{tabular}{lccc}
\toprule
Loci imposed & \# equations & Rank & Nullspace dim \\
\midrule
Leg 1 only     & 34  & 30 & 15 \\
Legs 1, 2      & 62  & 44 &  1 \\
Legs 1, 2, 3   & 83  & 44 &  1 \\
Legs 1, 2, 3, 4 & 99 & 44 &  1 \\
All 5 legs     & 115 & 44 &  1 \\
\bottomrule
\end{tabular}
\end{center}
In the $45$-parameter ansatz (all pairs of $10$ Mandelstam
invariants), two loci already suffice.
In the $15$-parameter ansatz of Part~I (planar variables only),
three loci are needed: the non-local survivor $\cN$ of
equation~\eqref{eq:after-leg2} persists after legs~$1$ and~$2$
and is killed only by leg~$3$.

\textbf{At $n = 6$ (numerical):}
\begin{center}
\small
\begin{tabular}{lcc}
\toprule
Loci imposed & Rank (of $84$) & Nullspace dim \\
\midrule
1 locus   & 51 & 33 \\
2 loci    & 76 &  8 \\
3 loci    & 81 &  3 \\
4 loci    & 82 &  2 \\
5 loci    & 83 &  1 \\
All 6     & 83 &  1 \\
\bottomrule
\end{tabular}
\end{center}
At $n = 6$, five loci are needed; leg~$6$ is the only
redundant one. No combination of three loci suffices
(all $\binom{6}{3} = 20$ triples give nullspace dimension
$\geq 3$). This \textbf{corrects} an earlier conjecture
that $\lfloor n/2 \rfloor$ loci suffice: at $n = 6$,
$\lfloor 6/2 \rfloor = 3$ is not enough; the correct
count is $n - 1 = 5$.
\end{remark}

\begin{remark}[Factorization structure]
\label{rmk:factorization}
At each propagator pole $X_{ij} = 0$, the residue of
$\cA_n^{\mathrm{tree}}$ factorizes as
$\cA_L^{\mathrm{tree}} \cdot \cA_R^{\mathrm{tree}}$,
where the diagonal $(i,j)$ splits the $n$-gon into an
$|L|$-gon and an $|R|$-gon ($|L| + |R| = n + 2$).
The number of terms in the residue equals
$C_{|L|-2} \cdot C_{|R|-2}$ (product of Catalan numbers),
matching the product of triangulation counts.
This is verified numerically at $n = 5, 6$ in
\texttt{computations/verify\_factorization.py}.
\end{remark}

% ---------------------------------------------------------------
\section{Conclusion}
\label{sec:conclusion}
% ---------------------------------------------------------------

We have proved the following sequence of results.

\begin{enumerate}[label=(\arabic*)]
  \item \textbf{$n = 5$ base case (Part~I).}
    The ansatz space $\cB_5$ is finite-dimensional;
    imposing the five cyclic $1$-zero conditions
    reduces it to a one-dimensional space spanned
    by $\cA_5^{\mathrm{tree}}$ (Theorem~\ref{thm:n5}).
    Verified symbolically by
    \texttt{computations/full\_ansatz\_n5.py}.

  \item \textbf{$n = 6$ verification (Part~II,
    Proposition~\ref{prop:n6}).}
    The $84$-dimensional ansatz $\cB_6$ is similarly
    reduced to a one-dimensional nullspace by the six
    cyclic $1$-zeros, verified numerically.
    The result holds also for the broader $165$-dimensional
    ansatz allowing higher-order poles.

  \item \textbf{General $n$ (Part~II, Theorem~\ref{thm:main}).}
    By strong induction, the same uniqueness result holds
    for all $n \geq 5$: cyclic $1$-zeros uniquely determine
    $\phi^3$ tree amplitudes. The induction uses the
    Homogeneity Lemma~\ref{lem:homogeneity},
    Dimensional Reduction (Proposition~\ref{prop:dimred}),
    and Inherited $1$-Zeros (Lemma~\ref{lem:inherited}).

  \item \textbf{Locality is derived, not assumed.}
    At each inductive step, finiteness of $B$ on the
    $1$-zero loci forces pole coefficients to zero,
    and the vanishing conditions kill the residual.
    Physical locality (poles only in planar channels)
    is a \emph{consequence} of the $1$-zero constraints
    together with the homogeneity of the ansatz space.

  \item \textbf{Independence of $1$-zero loci.}
    At $n = 5$, two or three loci suffice (depending on
    the ansatz). At $n = 6$, five of the six loci are
    needed. The earlier conjecture that $\lfloor n/2 \rfloor$
    loci suffice is falsified at $n = 6$; the correct
    count appears to be $n - 1$ (one locus is always
    redundant by cyclic symmetry).
\end{enumerate}

\subsection*{Remaining gap}

The proof of Theorem~\ref{thm:main} contains one flagged gap
in Step~4: the factorization argument for the BCFW residues
of a general $B \in \cB_n$ (not just $\cA_n^{\mathrm{tree}}$).
The claim that the residue at each BCFW pole factorizes as a
product of lower-point quantities, and hence vanishes when
$\widetilde{B}|_{\cZ_1} = 0$, relies on the combinatorial
correspondence between BCFW poles and sub-polygon factorisations
established in \cite[Sec.~3]{Rodina2024}. Transposing this
argument to the abstract ansatz setting $\cB_n$ requires
verifying that the pole structure of an arbitrary $B \in \cB_n$
admits the same sub-polygon decomposition as the Feynman-diagram
sum. At $n = 5$ and $n = 6$ this gap is bypassed by the
independent direct computation (Part~I and
Proposition~\ref{prop:n6}). For $n \geq 7$, either a direct
computation at $n = 7$ (feasible but computationally expensive)
or a self-contained algebraic proof of the factorization
structure would close this gap.

\subsection*{Open problems}

\begin{enumerate}[label=(\roman*)]
  \item Close the factorization gap in Step~4 by a
    self-contained algebraic argument, independent of
    \cite{Rodina2024}.

  \item Determine the sharp bound for the number of
    independent $1$-zero loci at general $n$.
    The data suggests $n - 1$; a proof would explain
    why exactly one locus is always redundant.

  \item Extend the framework to Yang--Mills and graviton
    amplitudes, where the ansatz space and $1$-zero
    structure are more intricate.
\end{enumerate}

\subsection*{Acknowledgements}

Computational verification was performed using
SymPy~\cite{SymPy2017} and NumPy; scripts are available in the
accompanying \texttt{computations/} directory.

\bigskip
\bibliography{references}
\bibliographystyle{unsrt}
